% Punto de partida (~1 min). Fuentes: EXP-000, EXP-001, ISS-001, ISS-002, ISS-005. Mismo orden que report_03_baseline.
\begin{frame}{Punto de partida y problemas detectados}
  \begin{itemize}
    \item Baseline: \texttt{Flatten} $\rightarrow$ ReLU $\rightarrow$ \texttt{Dense} sobre píxeles en $[0, 255]$; la ReLU es la identidad, así que es una regresión logística
    \item Fase 0: semilla fija y \emph{split} estratificado. Mean recall en validación: \textbf{\res{baseline-mean-recall}} (accuracy \res{baseline-val-acc})
    \item \textbf{Pérdida desbordada}: $\sim$\res{baseline-train-loss} frente a $\ln K \approx$ \res{chance-loss}; la accuracy oscila y hay clases que nunca se predicen
    \item \textbf{Callbacks casi inertes}: paciencias de \res{baseline-es-patience} y \res{baseline-rlr-patience} épocas en un entrenamiento de \res{baseline-epochs-cap}
    \item \textbf{Coste}: cada época releía los TIF; se cachean una vez
  \end{itemize}
  \vspace{0.2cm}
  \keymessage{Cada fase parte del ganador de la anterior; la validación elige cada hiperparámetro y el punto de parada.}
  \speakernote{1}{Por qué el modelo es lineal. La pérdida es enorme porque Adam mueve cada peso del orden de la lr en cada paso: con \res{baseline-inputs} entradas en $[0, 255]$, un paso desplaza un logit miles de unidades (primera época: \res{baseline-loss-first}). El run original del enunciado usaba otro split y no se compara.}
\end{frame}
